% !TEX root = /dev/null
\documentclass[11pt]{amsart}
\usepackage[localtheorems]{raeez-math-template}

\providecommand{\C}{\mathbb{C}}
\providecommand{\R}{\mathbb{R}}
\providecommand{\Z}{\mathbb{Z}}
\providecommand{\Q}{\mathbb{Q}}
\providecommand{\dbar}{\bar{\partial}}
\providecommand{\dbarstar}{\bar{\partial}^{\ast}}
\providecommand{\gfrak}{\mathfrak{g}}
\providecommand{\ad}{\operatorname{ad}}
\providecommand{\Tr}{\operatorname{Tr}}
\providecommand{\Sym}{\operatorname{Sym}}
\providecommand{\End}{\operatorname{End}}
\providecommand{\Rep}{\operatorname{Rep}}
\providecommand{\CoHA}{\operatorname{CoHA}}
\providecommand{\Obs}{\operatorname{Obs}}
\providecommand{\CE}{\operatorname{CE}}
\providecommand{\BV}{\operatorname{BV}}
\providecommand{\QME}{\operatorname{QME}}
\providecommand{\hCS}{\operatorname{hCS}}
\providecommand{\cA}{\mathcal{A}}
\providecommand{\cE}{\mathcal{E}}
\providecommand{\cO}{\mathcal{O}}
\providecommand{\cF}{\mathcal{F}}
\providecommand{\cH}{\mathcal{H}}
\providecommand{\cL}{\mathcal{L}}
\providecommand{\cU}{\mathcal{U}}
\providecommand{\Lat}{\Lambda}
\providecommand{\PhiFA}{\Phi^{\mathrm{FA}}}
\providecommand{\SpCh}{\operatorname{Sp}^{\mathrm{ch}}}
\providecommand{\Eone}{E_1}
\providecommand{\Etwo}{E_2}
\providecommand{\Ethree}{E_3}
\providecommand{\Yp}{Y^{+}}
\providecommand{\Woneinf}{\mathcal{W}_{1+\infty}}
\providecommand{\Heis}{\operatorname{Heis}}
\providecommand{\frakgDelta}{\mathfrak{g}_{\Delta_5}}
\providecommand{\kanom}{\kappa_{\operatorname{anom}}}
\providecommand{\kch}{\kappa_{\mathrm{ch}}}
\providecommand{\kchc}{\kappa_{\mathrm{ch}}^{\mathrm{compact}}}
\providecommand{\kchH}{\kappa_{\mathrm{ch}}^{\mathrm{Heis}}}
\providecommand{\kcat}{\kappa_{\mathrm{cat}}}
\providecommand{\kBKM}{\kappa_{\mathrm{BKM}}}
\providecommand{\kfib}{\kappa_{\mathrm{fiber}}}

\theoremstyle{plain}
\newtheorem{theorem}{Theorem}[section]
\newtheorem{proposition}[theorem]{Proposition}
\newtheorem{lemma}[theorem]{Lemma}
\newtheorem{corollary}[theorem]{Corollary}
\theoremstyle{definition}
\newtheorem{definition}[theorem]{Definition}
\newtheorem{construction}[theorem]{Construction}
\newtheorem{hypothesis}[theorem]{Hypothesis}
\theoremstyle{remark}
\newtheorem{remark}[theorem]{Remark}

\title{Holomorphic Chern-Simons Observables on $\C^3$ and $K3\times E$}
\author{Raeez Lorgat}
\date{2026-04-22}

\begin{document}
\maketitle

\begin{abstract}
Holomorphic Chern-Simons theory in complex dimension three has two
different one-loop slots.  The degree-one quantum master equation
obstruction is the quartic Chern-Weil class
\(I_4^\rho\in\Sym^4(\gfrak^\vee)^{\gfrak}\).  The quadratic Casimir
\(C_2^\rho\) contributes a degree-zero two-point counterterm and
renormalises the kinetic pairing.  Compact statements on
\(K3\times E\) therefore require an explicit BV package: orientation,
framing, gauge fixing, heat-kernel regulator, completed local
functionals, factorisation or TCFT sewing, and a trivialisation of the
class \(\kanom(K3\times E,\gfrak,\rho)\).  After this package is fixed,
the \(K3\)-fibre specialisation of the stage-one holomorphic
factorisation algebra gives an \(\Eone\)-chiral algebra on \(E\).  In
the abelian Gaussian sector it is the lattice Heisenberg algebra of the
chosen \(K3\)-lattice.  On the principal Hall-Borcherds recognition
locus, the rank-three hyperbolic Cartan shadow is transported to the
vertex envelope \(U_{\mathrm{ch}}(\frakgDelta)\), whose denominator is
the Borcherds product \(\Delta_5\) with
\(\kBKM(\Delta_5)=c_1(0)/2=10/2=5\).  The compact Hodge invariant
\(\kchc(K3\times E)=0\), the Heisenberg value \(\kchH=3\), the
Borcherds weight \(5\), and the fibre rank \(24\) come from four
separate constructions.
\end{abstract}

\tableofcontents

\section{BV datum and one-loop classes}
\label{sec:bv-datum}

\begin{construction}[Holomorphic Chern-Simons datum]
\label{cons:hcs-datum}
Let \(X\) be a complex threefold with a nowhere-vanishing holomorphic
volume form \(\Omega_X\).  Let \((\gfrak,\langle-,-\rangle)\) be a
finite-dimensional complex quadratic Lie algebra.  On an open
\(U\subset X\) the local BV fields are
\[
  \cE(U)=\Omega^{0,\bullet}(U,\gfrak)[1].
\]
The classical action is
\[
  S_{\mathrm{cl}}(\cA)=
  \int_U \Omega_X\wedge
  \left\langle
    \frac12\,\cA\wedge\dbar\cA+
    \frac16\,\cA\wedge[\cA,\cA],
    1
  \right\rangle ,
  \qquad
  \cA\in\Omega^{0,\bullet}(U,\gfrak).
\]
The shifted BV pairing is
\[
  \omega(\alpha,\beta)=
  \int_U\Omega_X\wedge\langle\alpha,\beta\rangle ,
\]
which pairs \(\Omega^{0,p}\) with compactly supported
\(\Omega^{0,3-p}\) and has BV degree \(-1\) after the shift.
\end{construction}

\begin{proposition}[Classical master equation]
\label{prop:cme}
The functional \(S_{\mathrm{cl}}\) satisfies
\[
  \{S_{\mathrm{cl}},S_{\mathrm{cl}}\}=0 .
\]
\end{proposition}

\begin{proof}
The Hamiltonian vector field of \(S_{\mathrm{cl}}\) is the
Chevalley differential of the Dolbeault dg Lie algebra
\[
  \bigl(\Omega^{0,\bullet}(U,\gfrak),\dbar,[-,-]_{\gfrak}\bigr).
\]
The equation follows from \(\dbar^2=0\), the derivation rule for
\(\dbar\), the Jacobi identity in \(\gfrak\), and invariance of the
pairing.
\end{proof}

\begin{definition}[Costello quantum datum]
\label{def:costello-quantum-datum}
A perturbative quantisation consists of a Costello gauge-fixing operator
\(Q^{\mathrm{GF}}\), heat kernels \(K_t\), regularised BV kernels
\(P_{\varepsilon,L}\), effective interactions
\[
  I[L]\in\cO_{\mathrm{loc}}(\cE_X)[[\hbar]],
\]
and a completed local functional complex on which
\[
  QI[L]+\frac12\{I[L],I[L]\}_{L}+\hbar\Delta_L I[L]=0
  \tag{\QME_L}
\]
is defined.  The obstruction at each quantum order lies in degree one in
this local BV complex.
\end{definition}

\begin{definition}[Quartic invariant]
\label{def:quartic-invariant}
For a representation \(\rho:\gfrak\to\End(V)\), set
\[
  I_4^\rho(x_1,x_2,x_3,x_4)
  =
  \frac1{4!}\sum_{\sigma\in S_4}
  \Tr_V\!\left(
    \rho(x_{\sigma(1)})\rho(x_{\sigma(2)})
    \rho(x_{\sigma(3)})\rho(x_{\sigma(4)})
  \right).
\]
Then \(I_4^\rho\in\Sym^4(\gfrak^\vee)^{\gfrak}\).  The adjoint
choice is denoted \(I_4^{\ad}\).
\end{definition}

\begin{proposition}[Quartic QME obstruction in complex dimension three]
\label{prop:quartic-qme-obstruction}
The nonabelian one-loop gauge obstruction for holomorphic
Chern-Simons theory on a complex threefold is represented in the local
Chevalley complex by
\[
  \Theta^\rho_U(\alpha_0,\alpha_1,\alpha_2,\alpha_3)
  =
  \int_U
  \left[
  I_4^\rho\bigl(
    \alpha_0,\partial\alpha_1,\partial\alpha_2,\partial\alpha_3
  \bigr)
  \right]_{(3,3)},
  \qquad
  \alpha_i\in\Omega^{0,\bullet}_c(U,\gfrak).
\]
Its obstruction class is
\[
  \kanom(X,\gfrak,\rho)
  =
  [\,\hbar\,\Theta^\rho\,]
  \in
  H^1_{\mathrm{loc}}
  \bigl(\Omega^{0,\bullet}(X,\gfrak)[1],\cO_{\mathrm{loc}}\bigr).
\]
For \(\rho=\ad\), this is the Costello-Li descent class represented in
ghost-field notation by \(\int_X\Tr_{\ad}(c\,F_A^3)\).
\end{proposition}

\begin{proof}
The relevant one-loop graph is the wheel with four external gauge
insertions.  Its Lie factor is the symmetrised trace of four Lie algebra
inputs, hence an element of \(\Sym^4(\gfrak^\vee)^{\gfrak}\).  The
holomorphic form factor carries three \(\partial\)-derivatives; in
complex dimension three this gives a holomorphic top density.  The
\((3,3)\)-component is a local functional of cohomological degree one.
Costello obstruction theory identifies this degree-one class as the
first obstruction to solving \((\QME_L)\).
\end{proof}

\begin{definition}[Quadratic Casimir coefficient]
\label{def:quadratic-casimir}
Let \(\{T_a\}\) be a basis of \(\gfrak\), and let \(\{T^a\}\) be the
dual basis for the chosen invariant pairing.  The representation
theoretic quadratic Casimir is
\[
  C_2^\rho=\sum_a\rho(T_a)\rho(T^a).
\]
For the adjoint representation of a simple Lie algebra, this is the
usual dual-Coxeter quadratic coefficient.
\end{definition}

\begin{proposition}[Two-point counterterm]
\label{prop:two-point-counterterm}
With the flat heat-kernel normalisation, the two-point graph contributes
a degree-zero local counterterm of the form
\[
  S_{\mathrm{ct}}^{(2)}
  =
  -\hbar\,\frac{C_2^\rho}{(4\pi)^3}
  \log(L/\varepsilon)
  \int_X\Omega_X\wedge\Tr_\rho(\cA\,\dbar\cA),
\]
up to the convention for the invariant pairing and the heat-kernel
normalisation.  This term renormalises the field and the kinetic
pairing.  The degree-one QME obstruction is the quartic class of
Proposition~\ref{prop:quartic-qme-obstruction}.
\end{proposition}

\begin{proof}
Closing two legs of the cubic vertex gives the colour contraction
\(\sum_a\rho(T_a)\rho(T^a)=C_2^\rho\).  The remaining two external
fields have the tensor type of the kinetic functional
\(\int\Omega_X\wedge\Tr_\rho(\cA\dbar\cA)\).  The short-distance
six-real-dimensional heat-kernel integral supplies the logarithmic
factor.  Its cohomological degree is zero, while
\(\Theta^\rho\) has degree one and arity four.
\end{proof}

\begin{corollary}[Real three-dimensional comparison]
\label{cor:three-dimensional-comparison}
The hCS BV complex has no secondary topological coupling of the real
three-dimensional theory.  Its one-loop data in complex dimension three
are the degree-one quartic class \(I_4^\rho\) and the degree-zero
two-point coefficient \(C_2^\rho\).
\end{corollary}

\section{Flat local model on \texorpdfstring{$\C^3$}{C3}}
\label{sec:flat-c3}

\begin{proposition}[Bochner-Martinelli propagator]
\label{prop:bm-propagator}
On \(\C^3\) with the Euclidean Hermitian metric and
\(Q^{\mathrm{GF}}=\dbarstar\), define
\[
  P_{\varepsilon,L}(z,w)
  =
  \int_{\varepsilon}^{L}(\dbarstar\otimes1)K_t(z,w)\,dt ,
  \qquad
  0<\varepsilon<L<\infty .
\]
Off the diagonal, the limit as \(\varepsilon\to0\) and
\(L\to\infty\) is the Bochner-Martinelli kernel
\[
  P_{\mathrm{BM}}(z,w)
  =
  \frac{2!}{(2\pi i)^3}
  \sum_{k=1}^{3}(-1)^{k-1}
  \frac{\overline{z_k-w_k}}{\|z-w\|^6}\,
  d\bar z_1\wedge\cdots\widehat{d\bar z_k}\cdots\wedge d\bar z_3 .
\]
The finite-cutoff kernel satisfies \(\dbar P_{\varepsilon,L}=K_L-K_\varepsilon\).
\end{proposition}

\begin{proof}
The scalar heat kernel is
\((4\pi t)^{-3}\exp(-\|z-w\|^2/(4t))\).  Applying
\(\dbarstar\) in one variable differentiates the Gaussian and
contributes the numerator \(\overline{z_k-w_k}\).  Integration in
\(t\) gives a smooth homotopy kernel at finite cutoffs, and the
telescoping identity for
\(\Delta_{\dbar}=\dbar\dbarstar+\dbarstar\dbar\) gives
\(\dbar P_{\varepsilon,L}=K_L-K_\varepsilon\).  Removing the cutoffs
off the diagonal recovers the displayed Bochner-Martinelli
representative.
\end{proof}

\begin{proposition}[Flat \(E_3\) PBW stalk]
\label{prop:flat-e3-pbw-stalk}
Let \(U\subset\C^3\) be a Stein polydisc, or let \(U=\C^3\) with
polynomial boundary conditions.  The classical hCS factorisation
algebra generated by
\[
  \mathscr L_{\hCS}(U)=\Omega^{0,\bullet}_c(U,\gfrak)
\]
is the \(E_3\)-holomorphic enveloping factorisation algebra.  Its
completed classical observables are
\[
  \Obs^{\mathrm{cl}}_{\hCS}(U)
  =
  C^\bullet_{\CE,\mathrm{cont}}
  \bigl(\mathscr L_{\hCS}(U)[1]\bigr)
  =
  \widehat{\Sym}\bigl(\mathscr L_{\hCS}(U)[1]^\vee[-1]\bigr).
\]
After a Costello quantum datum satisfying \((\QME_L)\) has been fixed,
the local quantum stalk is filtered by the \(E_3\) PBW filtration with
associated graded the completed signed symmetric algebra on
\(\mathscr L_{\hCS}[2]\).
\end{proposition}

\begin{proof}
Francis' \(E_n\) PBW theorem identifies the associated graded of the
\(E_n\)-enveloping filtration with the shifted symmetric algebra on the
generating Lie object.  At \(n=3\) the shift is \([2]\), the degree of
the Browder operation attached to
\[
  H^\bullet(\operatorname{Conf}_2(\R^3))\cong H^\bullet(S^2).
\]
The observable formula is the continuous Chevalley cochain algebra of
the local Dolbeault dg Lie algebra.  The quantum filtration exists only
after the scale-\(L\) QME has been solved in the completed BV complex.
\end{proof}

\begin{corollary}[Flat nonabelian scope]
\label{cor:flat-nonabelian-scope}
For abelian \(\gfrak\) the cubic interaction vanishes and the theory is
Gaussian.  For nonabelian \(\gfrak\), a flat \(\C^3\) quantisation
requires \(I_4^\rho=0\) or an explicit trivialisation of
\(\kanom(\C^3,\gfrak,\rho)\) in the chosen Costello local functional
complex.  A vanishing two-point bubble or a Casimir computation is
insufficient for the QME.
\end{corollary}

\begin{remark}[\(\C^3\) CoHA normalisation]
\label{rem:c3-coha-normalisation}
The toric cohomological Hall algebra attached to \(\C^3\) is
\[
  \CoHA(\C^3)\simeq \Yp(\widehat{\mathfrak{gl}}_1),
\]
the positive half of the affine Yangian in the Schiffmann-Vasserot
normalisation.  The full \(\Woneinf\) vertex algebra appears after the
appropriate Drinfeld double or centre, completion, and evaluation
representation have been constructed.
\end{remark}

\section{Compact product \texorpdfstring{$K3\times E$}{K3 x E}}
\label{sec:compact-k3e}

\begin{hypothesis}[Compact BV package]
\label{hyp:compact-bv-package}
A compact nonabelian holomorphic Chern-Simons quantisation theorem on a
Calabi-Yau threefold requires:
\begin{enumerate}
\item a compact Calabi-Yau threefold \((X,\Omega_X)\) with an
orientation and determinant-line framing compatible with the BV pairing;
\item a quadratic gauge Lie algebra or dg Lie algebra
\((\gfrak,\langle-,-\rangle)\), and a representation \(\rho\) whenever
traces are used, fixing \(I_4^\rho\) and \(C_2^\rho\);
\item an elliptic Costello gauge fixing, a heat-kernel regulator, and
an RG-compatible counterterm scheme;
\item a completed space of local functionals and observables on which
\(\Delta_L\), \(\{-,-\}_L\), and the RG flow are defined;
\item a factorisation or TCFT sewing structure with descent from
holomorphic polydiscs to \(X\);
\item a trivialisation of \(\kanom(X,\gfrak,\rho)\), by vanishing in
local BV cohomology, an explicit local counterterm, or an open-closed
Green-Schwarz term.
\end{enumerate}
\end{hypothesis}

\begin{theorem}[Conditional compact quantisation]
\label{thm:conditional-compact-quantisation}
Assume Hypothesis~\ref{hyp:compact-bv-package}.  If
\[
  \kanom(X,\gfrak,\rho)=0
\]
in the Costello local BV obstruction complex after the chosen
counterterms and closed-sector contributions are included, then the
effective interactions can be chosen to satisfy \((\QME_L)\) through
one loop.  Under the analogous higher-loop obstruction vanishing or
counterterm hypotheses, the completed observables form a quantum
factorisation algebra on \(X\).
\end{theorem}

\begin{proof}
Costello's obstruction theory solves the QME inductively in powers of
\(\hbar\).  At the first quantum step the obstruction is the degree-one
local class of Proposition~\ref{prop:quartic-qme-obstruction}, corrected
by the specified local or closed-sector counterterms.  Vanishing gives
the first quantum correction.  The same induction applies at later loop
orders inside the completed local functional complex.
\end{proof}

\begin{proposition}[Local and global data on \(K3\times E\)]
\label{prop:k3e-local-global}
Let \(X=K3\times E\), with
\[
  \Omega_X=\Omega_{K3}\wedge dz_E .
\]
At a point \(x\in X\), a holomorphic polydisc
\[
  D_x\simeq \C^2_{K3}\times\C_E
\]
has the same flat hCS local model as \(\C^3\).  The global compact
theory is governed by Hypothesis~\ref{hyp:compact-bv-package}; its
Dolbeault cohomology is
\[
  H^{0,\bullet}(K3\times E)
  =
  H^{0,\bullet}(K3)\otimes H^{0,\bullet}(E),
  \qquad
  (1,0,1)\otimes(1,1)=(1,1,1,1).
\]
Consequently
\[
  \kchc(K3\times E)
  =
  1-1+1-1
  =
  0,
  \qquad
  \kcat(K3\times E)=
  \chi(\cO_{K3})\chi(\cO_E)=2\cdot0=0.
\]
\end{proposition}

\begin{proof}
The local statement uses a holomorphic coordinate chart
\((z_1,z_2,w)\) with \((z_1,z_2)\) on \(K3\) and \(w\) on \(E\).  The
heat kernel has the usual short-time expansion
\[
  K_t^X(z,w)=K_t^{\C^3,\mathrm{loc}}(z,w)\bigl(1+t\,a_1(z,w)+O(t^2)\bigr),
\]
so the germ is the flat Dolbeault model.  Global observables retain the
compact harmonic modes and require the QME, regulator, completion, and
sewing data listed in Hypothesis~\ref{hyp:compact-bv-package}.  The
Hodge-number calculation is Kunneth:
\[
  H^{0,\bullet}(K3)=(1,0,1),\qquad
  H^{0,\bullet}(E)=(1,1).
\]
\end{proof}

\begin{remark}[Characteristic numbers and anomaly classes]
\label{rem:k3e-characteristic-numbers}
For \(K3\times E\),
\[
  \chi_{\mathrm{top}}(K3\times E)=24\cdot0=0.
\]
The second Chern class of the K3 factor belongs to the geometry of the
compact factor.  The open nonabelian hCS gauge obstruction is the local BV class
\(\kanom(K3\times E,\gfrak,\rho)\) determined by \(I_4^\rho\), the
trace normalisation, the regulator, and the descent representative.
The quadratic coefficient \(C_2^\rho\) remains a two-point kinetic
renormalisation datum.
\end{remark}

\section{The \texorpdfstring{$K3$}{K3} specialisation}
\label{sec:k3-specialisation}

\begin{definition}[Two-stage specialisation]
\label{def:two-stage-specialisation}
On the verified \(d=3\) loci, and after a witnessed admissible
specialisation datum \((\Sigma_2,C)\) has been fixed, the CY-to-chiral
assignment is
\[
  \Phi_3^{(\Sigma_2,C)}
  =
  \SpCh_{\Sigma_2,C}\circ\PhiFA_3 .
\]
Here \(\PhiFA_3\) is the stage-one \(E_3\)-holomorphic factorisation
object, and \(\SpCh_{\Sigma_2,C}\) is factorisation homology over
\(\Sigma_2\) followed by restriction to the curve \(C\).  For
\[
  X=K3\times E,\qquad
  \Sigma_2=K3,\qquad
  C=E,
\]
the specialised output is an \(\Eone\)-chiral algebra on \(E\).
\end{definition}

\begin{proposition}[\(E_1\) output in complex dimension three]
\label{prop:e1-output}
Let \(C\subset X\) be a holomorphic curve and assume the compact BV
package has been solved on a neighbourhood of \(C\).  Integration of
completed quantum observables in the normal directions gives
\[
  \mathcal A_C=
  \int_{N_{C/X}/C}\Obs^q_{\hCS}
  \in
  \operatorname{Alg}^{\mathrm{ch}}_{\Eone}(C).
\]
The \(\Etwo\)-structure belongs to the centre, double, or module layer,
for example
\[
  Z\bigl(\Rep^{\Eone}(\mathcal A_C)\bigr)
  \simeq
  \Rep^{\Etwo}\!\left(
    C^\bullet_{\CE,\mathrm{ch}}(\mathcal A_C,\mathcal A_C)
  \right)
\]
under the chiral BZFN hypotheses.
\end{proposition}

\begin{proof}
Stage one is \(E_3\)-holomorphic before specialisation.  Stage two
integrates over a complex surface and leaves one complex direction, so
factorisation homology drops the native operadic level to \(E_1\) on
the reference curve.  Braiding is recovered only after applying the
representation centre, Drinfeld double, or an equivalent module-layer
construction.
\end{proof}

\begin{construction}[Gaussian \(K3\)-pushforward]
\label{cons:gaussian-k3-pushforward}
Let \(L\subset H^{\mathrm{even}}(K3,\Z)\) be an even lattice with Mukai
pairing \((\, ,\,)_L\), and take the abelian coefficient algebra
\(\gfrak=L\otimes_\Z\C\).  For a small disc \(B\subset E\), the abelian
sector of the \(K3\)-pushforward is
\[
  \int_{K3}\Obs^q_{\hCS}(K3\times B,L\otimes\C)
  \simeq
  \widehat{\Sym}
  \bigl(L\otimes_\Z\Omega^{0,\bullet}_c(B)^\vee[-1]\bigr)[[\hbar]]_\star .
\]
The induced OPE is
\[
  \alpha(z)\beta(w)
  \sim
  \frac{(\alpha,\beta)_L}{(z-w)^2},
  \qquad
  \alpha,\beta\in L,
\]
and the resulting vertex algebra is the lattice Heisenberg algebra
\(\Heis_L\) on \(E\).
\end{construction}

\begin{proof}
The abelian BV theory is Gaussian.  Factorisation homology over the
compact \(K3\)-direction transfers the \(K3\) zero modes to coefficients
on \(E\), and the BV pairing contracts them by the Mukai form.  The
one-dimensional Green kernel on the elliptic disc supplies the double
pole, while the coefficient of the pole is the lattice pairing.  This is
the standard Heisenberg vertex algebra construction for an integral even
lattice.
\end{proof}

\begin{corollary}[Hyperbolic Cartan and Mukai fibre values]
\label{cor:heisenberg-values}
For the principal \(\Delta_5\) comparison, the relevant Cartan lattice is
the rank-three hyperbolic lattice \(\Lat^{2,1}_{II}\).  Hence
\[
  \kchH(K3\times E)=3 .
\]
For the full \(K3\) Mukai lattice
\[
  \widetilde{\Lat}_{K3}=H^{\mathrm{even}}(K3,\Z)
  \cong U^{\oplus4}\oplus E_8(-1)^{\oplus2}
\]
the rank is \(24\), giving
\[
  \kfib(K3\times E)=24 .
\]
The Heisenberg rank and the fibre rank record different specialisation
data.
\end{corollary}

\section{Borcherds recognition at the principal locus}
\label{sec:borcherds-recognition}

\begin{hypothesis}[Hall-Borcherds recognition datum]
\label{hyp:hall-borcherds-recognition}
The principal \(K3\times E\) Borcherds comparison consists of the
following data:
\begin{enumerate}
\item a compact Hall or reduced DT source for \(\operatorname{Perf}(K3\times E)\)
with orientation and compactness sufficient for Hall multiplication;
\item the \(K3\)-fibre specialisation
\(\SpCh_{K3,E}(\PhiFA_3(\operatorname{Perf}(K3\times E)))\);
\item a transport from this specialisation to the rank-three hyperbolic
lattice \(\Lat^{2,1}_{II}\);
\item imaginary-root multiplicities
\[
  \operatorname{mult}(n,l,m)=c_{\phi}(nm,l)
\]
from the primitive \(\Delta_5\) Borcherds input, with
\[
  c_{\phi}(0,0)=10;
\]
\item a PBW and no-extra-relations comparison identifying the transported
Hall presentation with the vertex envelope of \(\frakgDelta\).
\end{enumerate}
\end{hypothesis}

\begin{theorem}[Principal Borcherds output]
\label{thm:principal-borcherds-output}
Assume Hypothesis~\ref{hyp:compact-bv-package} and
Hypothesis~\ref{hyp:hall-borcherds-recognition} for \(X=K3\times E\).
Then the \(K3\)-specialised \(\Eone\)-chiral algebra on \(E\), after the
Hall-Borcherds transport, is
\[
  \SpCh_{K3,E}\bigl(\PhiFA_3(\operatorname{Perf}(K3\times E))\bigr)
  \simeq
  U_{\mathrm{ch}}(\frakgDelta).
\]
The denominator identity is
\[
  \frac{1}{64}\Delta_5(2Z)
  =
  e^{-2\pi i(\rho_3,z)}
  \prod_{\alpha\in\Delta_+}
  \bigl(1-e^{-2\pi i(\alpha,z)}\bigr)^{\operatorname{mult}(\alpha)},
\]
with
\[
  \rho_3=f_2-\frac12 f_3+f_{-2},
  \qquad
  \operatorname{mult}(n f_2+l f_3+m f_{-2})=c_{\phi}(nm,l).
\]
\end{theorem}

\begin{proof}
The recognition datum supplies a Hall presentation of the specialised
\(\Eone\)-chiral algebra and transports its Cartan to
\(\Lat^{2,1}_{II}\).  The PBW comparison identifies the real-root
generators with the real simple roots of the Borcherds presentation and
identifies the imaginary-root generators with the coefficients
\(c_{\phi}(nm,l)\).  Borcherds' automorphic product theorem then gives
the displayed denominator formula.  The factor \(64\) is the
normalisation used for the primitive \(\Delta_5(2Z)\) denominator.
\end{proof}

\begin{proposition}[Borcherds weight]
\label{prop:borcherds-weight}
In the primitive \(\Delta_5\) denominator normalisation,
\[
  \kBKM(\Delta_5)
  =
  \frac{c_1(0)}{2}
  =
  \frac{10}{2}
  =
  5.
\]
This is the automorphic weight of the Borcherds product.
\end{proposition}

\begin{proof}
Borcherds' singular-theta lift gives
\[
  \mathrm{wt}(\operatorname{BorLift}(\phi))=
  \frac{c_{\phi}(0,0)}2 .
\]
For the primitive \(\Delta_5\) input, the constant coefficient is
\(c_{\phi}(0,0)=10\).  Hence
\(\mathrm{wt}(\Delta_5)=5\), which is \(\kBKM(\Delta_5)\) in the
Vol III convention.
\end{proof}

\begin{corollary}[Four separate \(K3\times E\) invariants]
\label{cor:four-k3e-invariants}
The standard \(K3\times E\) values are
\[
  \kcat(K3\times E)=0,\qquad
  \kchc(K3\times E)=0,\qquad
  \kchH(K3\times E)=3,\qquad
  \kBKM(\Delta_5)=5,\qquad
  \kfib(K3\times E)=24.
\]
The categorical value is the Kunneth product
\(\chi(\cO_{K3})\chi(\cO_E)\).  The compact chiral value is the
Hodge supertrace of \(H^{0,\bullet}(K3\times E)\).  The Heisenberg value
is the rank of the \(\Lat^{2,1}_{II}\) Cartan shadow.  The Borcherds
value is \(c_1(0)/2\).  The fibre value is the rank of the full Mukai
lattice.  Any equality among these numbers requires a theorem for the
named construction.
\end{corollary}

\section{Automorphisms and equivariant refinements}
\label{sec:omega-scope}

\begin{proposition}[Equivariant scope of \(K3\times E\)]
\label{prop:equivariant-scope}
The toric \(\Omega\)-background on \(\C^3\) uses a dense
\((\C^\ast)^3\)-action with isolated fixed points.  For a compact K3
surface,
\[
  H^0(K3,T_{K3})=0,
\]
so the identity component of the automorphism group is trivial.  For a
generic product \(K3\times E\),
\[
  \operatorname{Aut}^0(K3\times E)=E,
\]
acting by translations on the elliptic factor.  Thus the compact product
supplies at most the elliptic spectral parameter.  Finite Nikulin-Mukai
symplectic automorphisms of \(K3\) give discrete CHL twistings; continuous
toric deformation parameters require a torus action.
\end{proposition}

\begin{proof}
The holomorphic symplectic form identifies \(T_{K3}\) with
\(\Omega^1_{K3}\), and \(H^0(K3,\Omega^1_{K3})=0\).  Hence a compact
K3 has no nonzero holomorphic vector field.  The elliptic curve
contributes its translation group.  Nikulin and Mukai classify finite
symplectic automorphism groups of K3 surfaces; these groups provide
finite equivariant sectors in place of a continuous torus action on the
K3 factor.
\end{proof}

\begin{remark}[Toric and compact outputs]
\label{rem:toric-compact-outputs}
The \(\C^3\) toric output
\(\CoHA(\C^3)\simeq\Yp(\widehat{\mathfrak{gl}}_1)\) belongs to the
affine fixed-point calculation.  The \(K3\times E\) compact output uses
the \(K3\)-fibre specialisation, the compact BV package, and the
Hall-Borcherds recognition datum.  These are different geometric
inputs: flat affine \(\C^3\), compact \(K3\times E\), and transverse
\(K3\)-integration.
\end{remark}

\section{Conclusion}
\label{sec:conclusion}

The local hCS theory on \(\C^3\) supplies the Bochner-Martinelli
propagator and the \(E_3\) PBW stalk after the QME has been solved.
The nonabelian QME obstruction in complex dimension three is the
quartic class \(I_4^\rho\); \(C_2^\rho\) is a two-point kinetic
coefficient.  The compact theory on \(K3\times E\) is a separate
problem requiring anomaly trivialisation, orientation, framing,
gauge-fixing, regulator, completion, and sewing.  Under those
hypotheses, the \(K3\)-fibre specialisation gives an \(\Eone\)-chiral
algebra on \(E\).  The abelian sector is the lattice Heisenberg algebra
of the chosen \(K3\)-lattice.  On the principal Hall-Borcherds locus,
the transported rank-three Cartan and the \(\phi_{0,1}\) multiplicities
recognise \(U_{\mathrm{ch}}(\frakgDelta)\), with
\(\kBKM(\Delta_5)=10/2=5\).  The values
\[
  0,\quad 0,\quad 3,\quad 5,\quad 24
\]
record \(\kcat\), \(\kchc\), \(\kchH\), \(\kBKM\), and \(\kfib\),
respectively.

\begin{thebibliography}{99}
\bibitem{Lorgat2020}
R. Lorgat, ``A Borcherds lift of the weak Jacobi form \(\phi_{0,1}\),
generalised Borcherds-Kac-Moody superalgebras and the Igusa cusp form
\(\Delta_5\)'', April 2020.

\bibitem{GritsenkoNikulin1995}
V. Gritsenko and V. Nikulin, ``Siegel automorphic form corrections of
some Lorentzian Kac-Moody Lie algebras'', \texttt{alg-geom/9504006}
(1995).

\bibitem{Borcherds1998}
R. Borcherds, ``Automorphic forms with singularities on
Grassmannians'', \texttt{alg-geom/9609022} (1996),
\emph{Invent. Math.} 132 (1998), 491-562.

\bibitem{Costello2011}
K. Costello, \emph{Renormalization and Effective Field Theory}, AMS
Mathematical Surveys 170 (2011).

\bibitem{CostelloGwilliam2017}
K. Costello and O. Gwilliam, \emph{Factorization Algebras in Quantum
Field Theory}, Volume 1, Cambridge (2017).

\bibitem{CostelloGwilliam2021}
K. Costello and O. Gwilliam, \emph{Factorization Algebras in Quantum
Field Theory}, Volume 2, Cambridge (2021).

\bibitem{CostelloLi2020}
K. Costello and S. Li, ``Quantization of open-closed BCOV theory I'',
\texttt{arXiv:1505.06703}, published 2020.

\bibitem{Nikulin1980}
V. Nikulin, ``Finite groups of automorphisms of Kahlerian K3 surfaces'',
\emph{Trudy Moskov. Mat. Obshch.} 38 (1980).

\bibitem{Mukai1988}
S. Mukai, ``Finite groups of automorphisms of K3 surfaces and the
Mathieu group'', \emph{Invent. Math.} 94 (1988), 183-221.

\bibitem{SchiffmannVasserot2013}
O. Schiffmann and E. Vasserot, ``The elliptic Hall algebra and the
\(K\)-theory of the Hilbert scheme of \(\mathbb A^2\)'',
\emph{Duke Math. J.} 162 (2013), 279-366.

\bibitem{MaulikOkounkov2012}
D. Maulik and A. Okounkov, ``Quantum groups and quantum cohomology'',
\texttt{arXiv:1211.1287}.

\bibitem{FLM1988}
I. Frenkel, J. Lepowsky and A. Meurman,
\emph{Vertex Operator Algebras and the Monster}, Academic (1988).

\bibitem{LurieHA}
J. Lurie, \emph{Higher Algebra}.
\end{thebibliography}

\end{document}
